\section{Materiais e Métodos}

\subsection{Materiais e Reagentes}

Para a realização do experimento, foram utilizados:
\begin{itemize}
    \item dois balões volumétricos de $50\un{mL}$ (identificados como 1 e 2);
    \item béqueres de $50$, $100$ e $150\un{mL}$;
    \item duas buretas, suporte universal e garras;
    \item bastão de vidro, pipeta Pasteur, cronômetro e termômetro;
    \item banho de gelo;
    \item Solução 1: $4{,}28\un{g}$ de \ce{KIO3} em $1\un{L}$ de solução;
    \item Solução 2: $0{,}852\un{g}$ de \ce{NaHSO3}, $4\un{mL}$ de \ce{H2SO4} concentrado e $50\un{mL}$ de suspensão de amido em $1\un{L}$ de solução;
    \item solução de iodato de concentração desconhecida;
    \item água destilada.
\end{itemize}

\subsection{Métodos}

O procedimento experimental foi dividido em duas partes, conforme a apostila da disciplina\cita{1}:
\begin{itemize}
    \item influência da concentração na velocidade da reação e determinação da concentração de uma solução desconhecida (parte A);
    \item influência da temperatura na velocidade da reação (parte B).
\end{itemize}

\subsubsection{Parte A -- Influência da concentração}

As buretas foram condicionadas com pequenas porções das respectivas soluções (bureta 1 com a Solução 1 e bureta 2 com a Solução 2), fixadas no suporte e completadas até a marca zero. Para cada amostra (A a E), transferiu-se diretamente da bureta para o balão volumétrico 1 o volume de Solução 1 indicado na Tabela~\ref{tab:composicao}, completando-se com água destilada até a marca de $50\un{mL}$ (ajuste final do menisco com a pipeta Pasteur), e agitou-se. O mesmo foi feito no balão 2 com $10{,}00\un{mL}$ da Solução 2.

\begin{table}[H]
\centering
\caption{Volumes das soluções utilizados no preparo de cada amostra.}
\label{tab:composicao}
\begin{tabular}{ccc}
\toprule
Amostra & Solução 1 (mL) & Solução 2 (mL) \\
\midrule
A & 5,00 & 10,00 \\
B & 10,00 & 10,00 \\
C & 15,00 & 10,00 \\
D & 20,00 & 10,00 \\
E & 25,00 & 10,00 \\
Desconhecida & -- & 10,00 \\
\bottomrule
\end{tabular}
\fonte{Adaptada da apostila da disciplina\cita{1}.}
\end{table}

O conteúdo do balão 1 foi transferido para um béquer de $150\un{mL}$, contendo um bastão de vidro, e o do balão 2 para um béquer de $100\un{mL}$. A solução 2 foi então vertida sobre a solução 1, disparando-se o cronômetro imediatamente; a mistura foi agitada por 2 a 3 segundos e anotou-se o tempo necessário para o aparecimento da cor azul. Os béqueres foram lavados com água destilada entre os ensaios. Por fim, o procedimento foi repetido usando a solução desconhecida no lugar da Solução 1. Com os dados, construiu-se o gráfico do inverso do tempo em função da concentração de \ce{KIO3} após a mistura, ajustou-se uma reta pelo método dos mínimos quadrados (Seção~\ref{sec:mmq}) e, a partir dela, determinou-se a concentração da solução desconhecida.

\subsubsection{Parte B -- Influência da temperatura}

Prepararam-se, como na parte A, as soluções correspondentes à amostra C: $15{,}0\un{mL}$ da Solução 1 no balão 1 e $10{,}0\un{mL}$ da Solução 2 no balão 2, ambas completadas até $50\un{mL}$. As soluções foram transferidas para dois béqueres, que foram colocados em banho de gelo, com um termômetro no meio de cada solução, até atingirem $15\un{^\circ C}$. A solução 2 foi então vertida na solução 1, fora do banho, e o tempo até o aparecimento da cor azul foi cronometrado. O procedimento foi repetido com as soluções resfriadas a uma temperatura entre $4$ e $6\un{^\circ C}$. Os resultados foram comparados com o da amostra C da parte A, realizada à temperatura ambiente, que foi de $25\un{^\circ C}$.

\subsection{Fundamentação Teórica}

\subsubsection{Reações envolvidas}

As reações que ocorrem no relógio de iodo são\cita{1}:
\begin{gather}
\ce{IO3- + H+ <=> HIO3} \label{eq:r1}\\
\ce{2 HIO3 + 5 H2SO3 -> I2 + 5 H2SO4 + H2O} \label{eq:r2}\\
\ce{I2 + H2SO3 + H2O -> 2 HI + H2SO4} \label{eq:r3}
\end{gather}
em que o ácido sulfuroso provém do hidrogenossulfito em meio ácido ($\ce{HSO3- + H+ <=> H2SO3}$). A reação \eqref{eq:r2} é lenta, enquanto a \eqref{eq:r3} é muito rápida. Assim, todo iodo formado é imediatamente consumido enquanto houver \ce{H2SO3}. Quando este se esgota, o iodo se acumula e forma com o amido um complexo de cor azul intensa, marcando o final da reação.

\subsubsection{Concentrações após a mistura}

As concentrações das soluções estoque são obtidas de $c=m/(M\,V)$. Com as massas molares $M(\ce{KIO3})=214{,}00\un{g/mol}$ e $M(\ce{NaHSO3})=104{,}06\un{g/mol}$:
\begin{gather*}
c_1=\frac{4{,}28}{214{,}00\cdot1{,}000}=2{,}000\times10^{-2}\un{mol/L},\\
c_2=\frac{0{,}852}{104{,}06\cdot1{,}000}=8{,}188\times10^{-3}\un{mol/L}.
\end{gather*}
Como a quantidade de matéria se conserva na diluição ($c_iV_i=c_fV_f$) e o volume final da mistura é $V_f=50+50=100\un{mL}$, as concentrações após a mistura são
\begin{equation}
[\ce{KIO3}]=\frac{c_1V_1}{V_f},
\qquad
[\ce{NaHSO3}]=\frac{c_2\cdot10{,}00}{100}=8{,}19\times10^{-4}\un{mol/L},
\label{eq:diluicao}
\end{equation}
em que $V_1$ é o volume de Solução 1 utilizado. A concentração de \ce{NaHSO3} é a mesma em todas as amostras.

\subsubsection{Velocidade de reação e tempo de reação}

A velocidade média de consumo do hidrogenossulfito no intervalo de $0$ a $t$ é
\[
\bar{v}=\frac{[\ce{HSO3-}]_0-[\ce{HSO3-}]_t}{t}=\frac{[\ce{HSO3-}]_0}{t},
\]
pois, no instante em que surge a cor azul, o hidrogenossulfito foi totalmente consumido. Como $[\ce{HSO3-}]_0$ é igual em todos os ensaios, tem-se $\bar{v}\propto 1/t$: o inverso do tempo (tempo recíproco) é proporcional à velocidade média da reação\cita{1}.

Supondo uma lei de velocidade da forma $v=k\,[\ce{IO3-}]^n[\ce{HSO3-}]^m[\ce{H+}]^p$, com $[\ce{HSO3-}]_0$ e $[\ce{H+}]$ fixos, obtém-se
\begin{equation}
\frac{1}{t}=K\,[\ce{KIO3}]^n
\quad\Longrightarrow\quad
\ln\frac{1}{t}=\ln K+n\ln[\ce{KIO3}],
\label{eq:ordem}
\end{equation}
em que $K$ reúne as constantes. Se a reação for de primeira ordem em relação ao iodato ($n=1$), $1/t$ é diretamente proporcional a $[\ce{KIO3}]$, o que justifica o uso de uma reta como curva de calibração. A ordem $n$ pode ser obtida pelo coeficiente angular do gráfico de $\ln(1/t)$ em função de $\ln[\ce{KIO3}]$.

\subsubsection{Efeito da temperatura -- equação de Arrhenius}

A dependência da constante de velocidade com a temperatura é descrita pela equação de Arrhenius\cita{2,3},
\[
k=A\,e^{-E_a/RT},
\]
em que $A$ é o fator pré-exponencial, $E_a$ a energia de ativação, $R=8{,}314\un{J\,mol^{-1}\,K^{-1}}$ e $T$ a temperatura absoluta. Como na parte B as concentrações são iguais em todos os ensaios, $1/t\propto k$, e
\begin{equation}
\ln\frac{1}{t}=C-\frac{E_a}{R}\cdot\frac{1}{T}.
\label{eq:arrhenius}
\end{equation}
Assim, o gráfico de $\ln(1/t)$ em função de $1/T$ é uma reta de coeficiente angular $-E_a/R$. Para dois ensaios, a Equação~\eqref{eq:arrhenius} fornece diretamente
\begin{equation}
E_a=\frac{R\,\ln(t_2/t_1)}{1/T_2-1/T_1},
\label{eq:Ea2}
\end{equation}
sendo $t_1$ e $t_2$ os tempos de reação às temperaturas $T_1>T_2$.

\subsubsection{Método dos mínimos quadrados e incertezas}
\label{sec:mmq}

Para ajustar uma reta $y=a+bx$ a $N$ pontos $(x_i,y_i)$, minimiza-se $S=\sum_i(y_i-a-bx_i)^2$, o que resulta em
\[
b=\frac{\sum_i(x_i-\bar{x})(y_i-\bar{y})}{S_{xx}},
\qquad
a=\bar{y}-b\,\bar{x},
\qquad
S_{xx}=\sum_i(x_i-\bar{x})^2,
\]
com incertezas
\[
u(b)=\frac{s_r}{\sqrt{S_{xx}}},
\qquad
u(a)=s_r\sqrt{\frac{1}{N}+\frac{\bar{x}^2}{S_{xx}}},
\qquad
s_r=\sqrt{\frac{\sum_i(y_i-a-bx_i)^2}{N-2}}.
\]
A qualidade do ajuste é avaliada pelo coeficiente de correlação $r$. Para uma amostra cuja resposta medida é $y_x$, a concentração interpolada na curva de calibração é $x_x=(y_x-a)/b$, com incerteza\cita{4}
\begin{equation}
u(x_x)=\frac{s_r}{b}\sqrt{1+\frac{1}{N}+\frac{(y_x-\bar{y})^2}{b^2S_{xx}}}.
\label{eq:ucal}
\end{equation}
